\chapter{Modelos}
\label{chap:modelos}

Este capítulo describe las dos familias de modelos que definen el TFG. La evidencia empírica cerrada corresponde a la CNN y a sus variantes de adaptación de dominio. Los modelos de visión y lenguaje no se presentan como adorno futuro, sino como el segundo brazo de la comparación metodológica: una forma alternativa de usar los mismos pocos ejemplos sin actualizar pesos.

\section{Entrada común}

Todos los modelos trabajan con imágenes de ECG. En la CNN, cada imagen corresponde a una derivación precordial derecha renderizada y redimensionada a 224 por 224 píxeles. En la línea VLM prevista, la misma imagen se incluirá en un mensaje multimodal junto con instrucciones y, si procede, ejemplos en contexto. La comparación solo tiene sentido porque ambos enfoques reciben la misma evidencia visual y se evalúan por paciente.

\begin{figure}[H]
  \centering
  \safeincludegraphics[width=0.96\textwidth]{results/fig_pipeline_experimental.pdf}
  \caption{Recorridos de entrada y evaluación: CNN cerrada y VLM/ICL preparado como comparador de pocos datos.}
  \label{fig:model_routes}
\end{figure}

\section{Dos formas de usar \texorpdfstring{\(K\)}{K}}

El presupuesto \(K\) tiene significados distintos según el modelo. En la CNN, \(K\) determina cuántos pacientes etiquetados se usan para optimizar los pesos. El aprendizaje ocurre durante el entrenamiento y la inferencia posterior solo aplica la red ajustada. En un VLM con aprendizaje en contexto, \(K\) determina cuántas demostraciones se incluyen en el prompt. Los pesos permanecen fijos y el aprendizaje se produce de forma implícita durante la generación.

Esta diferencia es la razón de la comparación. Si \(K\) es muy pequeño, entrenar una CNN puede producir alta varianza, sobreajuste o dependencia de atajos visuales. Un VLM, en cambio, podría aprovechar conocimiento visual previo y usar los ejemplos para fijar el formato de la tarea. También puede fallar: puede ignorar la imagen, apoyarse en el texto o devolver respuestas inválidas. Por eso la rama VLM requiere controles específicos.

\section{CNN ResNet18}

El modelo supervisado ejecutado es una ResNet18 de \texttt{torchvision} con pesos iniciales \texttt{DEFAULT}. La cabeza de clasificación se sustituye por una salida binaria. En los experimentos sintéticos, la etiqueta positiva indica que el paciente simulado cumple simultáneamente los hallazgos operativos asociados al patrón Brugada derivado. En HUCA, la etiqueta positiva indica caso Brugada confirmado frente a control.

La pérdida es \texttt{BCEWithLogitsLoss}. Cuando el pliegue lo requiere, la clase positiva se pondera automáticamente a partir de los ejemplos de entrenamiento. El optimizador es AdamW, con tasa de aprendizaje \(3\times10^{-4}\), lote de 32 imágenes y 20 épocas máximas. La selección de umbral se realiza en validación interna mediante exactitud equilibrada derivada por paciente; si no hay validación suficiente, se conserva el umbral 0,5.

\begin{table}[H]
  \centering
  \caption{Configuración CNN utilizada en las campañas finales.}
  \label{tab:modelo_cnn_final}
  \small
  \begin{tabularx}{0.92\textwidth}{p{4cm}X}
    \toprule
    Propiedad & Valor \\
    \midrule
    Arquitectura & ResNet18 con cabeza binaria \\
    Inicialización & Pesos \texttt{DEFAULT} de Torchvision para ResNet18 \\
    Entrada & RGB, 224 por 224 píxeles \\
    Pérdida & \texttt{BCEWithLogitsLoss}, \texttt{loss\_pos\_weight=auto} \\
    Optimizador & AdamW \\
    Tasa de aprendizaje & \(3\times10^{-4}\) \\
    Épocas máximas & 20 \\
    Tamaño de lote & 32 \\
    Presupuestos & \(K=\{2,4,8,16,32\}\) \\
    Semillas & 42, 123 y 2026 \\
    Umbral & Seleccionado por validación con exactitud equilibrada por paciente \\
    Agregación & Media de probabilidades de V1, V2 y V3 antes del umbral \\
    \bottomrule
  \end{tabularx}
\end{table}

\section{Validación por paciente}

El entrenamiento y la evaluación se organizan mediante validación leave-one-out. En cada pliegue se reserva un paciente completo y se entrena con un subconjunto de \(K\) pacientes del resto. Esto evita que imágenes de un mismo paciente aparezcan simultáneamente en entrenamiento y test. En HUCA esta restricción es indispensable porque las tres derivaciones de un paciente comparten adquisición, calibración y contexto clínico.

El informe de cada ejecución guarda:

\begin{itemize}
  \item configuración de entrenamiento;
  \item pacientes seleccionados como ejemplos;
  \item probabilidades por imagen;
  \item predicción agregada por paciente;
  \item umbral elegido;
  \item matriz de confusión y métricas;
  \item paneles Grad CAM para inspección cualitativa.
\end{itemize}

\section{Grad CAM}

Grad CAM se utiliza como diagnóstico cualitativo, no como explicación clínica suficiente \cite{selvaraju2017}. El mapa indica regiones que influyen en el logit del modelo, pero no prueba que la red haya aprendido un criterio médico correcto. En esta memoria se usa para comprobar si las activaciones se concentran en zonas plausibles del trazado y para detectar errores en los que una atención visual razonable termina en una decisión equivocada.

\begin{figure}[H]
  \centering
  \safeincludegraphics[width=0.92\textwidth]{results/cnn_simulator_qrs_gradcam_example.png}
  \caption{Ejemplo Grad CAM del modelo CNN sobre el conjunto sintético.}
  \label{fig:gradcam_sintetico_modelos}
\end{figure}

\section{Adaptación de dominio}

La adaptación de dominio se entrena con el conjunto sintético como origen etiquetado y HUCA como destino no etiquetado durante el ajuste. La evaluación final sigue realizándose sobre HUCA por paciente. Se revisan cuatro variantes:

\begin{itemize}
  \item \textbf{SSL}. Preentrenamiento SimCLR durante tres épocas sobre imágenes del dominio destino y posterior ajuste supervisado. En el código se reproduce con \texttt{METHOD=ssl}, que internamente llama al pipeline con \texttt{domain\_adaptation=none} y \texttt{ssl\_pretrain\_epochs=3}.
  \item \textbf{CORAL}. Penalización de covarianzas para aproximar la estructura de características de origen y destino.
  \item \textbf{MMD}. Penalización de discrepancia de medias en kernel con escalas 0,5, 1,0, 2,0 y 4,0.
  \item \textbf{DANN}. Clasificador de dominio adversario con inversión de gradiente \cite{ganin2016}.
\end{itemize}

Las pérdidas de adaptación se ponderan con 0,1 y se activan de forma progresiva durante el entrenamiento. Esta rama responde a una pregunta concreta: si la CNN falla por cambio de dominio visual, alinear representaciones debería recuperar parte del rendimiento. Los resultados muestran que la mejora existe, pero es demasiado pequeña para cerrar el problema.

\section{Modelos de visión y lenguaje}

La línea VLM queda preparada para la siguiente fase experimental. Su objetivo será comparar aprendizaje supervisado con aprendizaje en contexto bajo el mismo presupuesto \(K\). La CNN modifica pesos con ejemplos etiquetados; el VLM mantiene pesos fijos y recibe ejemplos dentro del mensaje. Esta diferencia convierte a los VLM en una alternativa especialmente relevante cuando etiquetar nuevos ECG es lento o caro.

El modelo abierto principal previsto es Gemma 4 y la referencia médica prevista es MedGemma. Ambos se servirán localmente mediante vLLM cuando haya recursos disponibles. La respuesta esperada será un JSON validado, sin explicación libre obligatoria, para evitar que una salida con formato incorrecto se mezcle con una predicción válida.

\begin{figure}[H]
  \centering
  \safeincludegraphics[width=0.96\textwidth]{results/fig_vlm_prompt_pipeline.pdf}
  \caption{Canalización VLM/ICL prevista: selección de ejemplos, prompt multimodal, modelo local, validación JSON y controles.}
  \label{fig:vlm_prompt_pipeline}
\end{figure}

\section{Controles previstos para VLM}

Cuando se ejecute la campaña VLM, deberán conservarse al menos tres condiciones:

\begin{itemize}
  \item \textbf{Normal}: ejemplos con imagen y respuesta correcta.
  \item \textbf{Respuestas permutadas}: las imágenes de ejemplo se mantienen, pero las respuestas se barajan para medir dependencia real de la correspondencia imagen-etiqueta.
  \item \textbf{Ejemplos sin imagen}: el contexto conserva texto y respuestas, pero elimina la evidencia visual de las demostraciones.
\end{itemize}

Estos controles son necesarios porque un VLM puede aprovechar sesgos textuales, frecuencia de etiquetas o formato del mensaje sin interpretar la morfología del ECG. Si el rendimiento no cae con respuestas permutadas o sin imágenes, la mejora no podría atribuirse a aprendizaje multimodal real. Hasta que no existan esos resultados, el VLM se mantiene como procedimiento reproducible pendiente.

\section{Criterio de comparación}

La comparación principal entre métodos se hará con exactitud equilibrada por paciente. F1, sensibilidad, especificidad, área ROC y precisión media se reportan como métricas complementarias. También se compararán coste de entrenamiento, coste por inferencia, tasa de respuestas inválidas, latencia y necesidad de hardware.

La lectura esperada no será "CNN contra VLM" como competición abstracta. Será una curva de decisión: con qué \(K\), qué coste y qué robustez conviene entrenar un modelo específico; y en qué zona de pocos datos puede rentar usar ICL para obtener una señal preliminar de triaje. En el cierre actual, todas las conclusiones cuantitativas proceden de CNN y adaptación de dominio. Esa base es precisamente el punto de partida de la comparación VLM.
